\section{Hope: the opportunity to save lives if we act}
\label{sec:hope}

Not every successful health law is built on fear. The most durable ones also
offer hope: earlier cancer detection, reduced physician burnout, better rural
access, and faster discovery. For medical AI, the hopeful claim is concrete. We
have an opportunity to save lives if we act, and the tools to do it responsibly
already exist and have been tested.

[DRAFTING INSTRUCTIONS: in the full-narrative, ground the hope appeal in
\texttt{review/references/narrative\_refs.bib} entry \texttt{topol2019high}
(human and artificial intelligence converging in medicine) and in the author's
simulation lineage in \texttt{review/references/author\_works.bib}
(\texttt{Kawchak2026AcceleratedPatientPredictionPhysical},
\texttt{Kawchak2025AcceleratingFDAComplianceCost}). Reproduce Mermaid figure
\texttt{review/mermaid/diagrams/06-evidence-to-law-lineage.md} as colored TikZ to
show that the hope rests on documented work. Add a body-width table of four
near-term benefits and the bill mechanism that enables each.]

\begin{narrativefig}
\node[nftwo] (a) at (0,0) {Verified tools};
\node[nfdec] (b) at (3.8,0) {Gates pass};
\node[nfhope] (c) at (7.6,0) {Lives saved};
\draw[nfedge] (a)--(b); \draw[nfedge] (b)--(c);
\end{narrativefig}
\vspace{-0.2cm}
\figcaption{Figure 5. Hope grounded in verified tools (placeholder; expanded from
Mermaid 06 in the full-narrative).}

Hope, paired with evidence, gives a legislator something to vote for rather than
only something to vote against.
